\begin{filecontents*}{\jobname.bib}
@article{hosseini2020,
  author  = {Hosseini, Seyed},
  title   = {Single Author Study on Software Verification},
  journal = {Journal of Systems},
  year    = {2020},
  volume  = {10},
  pages   = {100--110}
}
@article{du2021,
  author  = {Du, Wei and Gu, Chen},
  title   = {Collaborative Systems with Two Authors},
  journal = {Communications of ACM},
  year    = {2021},
  volume  = {20},
  pages   = {200--215}
}
@article{smith2022,
  author  = {Smith, John and Jones, Alice and Brown, Bob},
  title   = {Scalable Algorithms for Modern Networks},
  journal = {IEEE Transactions on Computers},
  year    = {2022},
  volume  = {30},
  pages   = {300--320}
}
@article{lee2023,
  author  = {Lee, Min and others},
  title   = {Large-Scale Experimental Evaluation},
  journal = {ACM Computing Surveys},
  year    = {2023},
  volume  = {40},
  pages   = {400--420}
}
@article{dimonda2025,
  author  = {Di Monda, Marco},
  title   = {On Compound Surnames in Citation Systems},
  journal = {Journal of Documentation},
  year    = {2025},
  volume  = {50},
  pages   = {500--515}
}
@manual{manual2024,
  title        = {Antigravity System Manual},
  organization = {Open Tools Project},
  year         = {2024}
}
@article{yamada2020,
  author  = {山田 太郎},
  title   = {日本語単著論文の引用表記法},
  journal = {情報処理学会論文誌},
  year    = {2020},
  volume  = {61},
  number  = {1},
  pages   = {1--8}
}
@article{sato2021,
  author  = {佐藤 一郎 and 鈴木 二郎},
  title   = {共著論文における連記表現の分析},
  journal = {情報処理学会論文誌},
  year    = {2021},
  volume  = {62},
  number  = {2},
  pages   = {9--16}
}
\end{filecontents*}

\documentclass[uplatex]{jsarticle}
\usepackage[dvipdfmx]{graphicx}
\usepackage{authcite}
\usepackage{url}

\title{\textsf{authcite} パッケージ\\取扱説明書兼サンプル文書}
\author{\texttt{authcite} 開発チーム}
\date{2026年10月29日 (v1.0)}

\begin{document}
\maketitle

\begin{abstract}
\textsf{authcite} は，\LaTeX において文献参照コマンド \verb|\cite| の直前に著者名（姓）を自動付与するパッケージです．
手動で「〇〇ら」と入力する手間を省き，文献データベース（Bib\TeX）から著者数や姓を正確に抽出して整形します．
日本語（「〜ら」「〜と〜」）および英語（「\textit{et al.}」「and」）の両方に対応し，2人共著の表記（省略または連記）も柔軟に切り替え可能です．
\end{abstract}

\section{はじめに}
学術論文や技術文書では，文献番号だけでなく「Hosseini[1]によれば」「Duら[2]の研究では」のように，
筆頭著者名を添えて引用するスタイルが広く用いられます．
しかし，著者の人数（1名，2名，3名以上）や表記規則（英語の \textit{et al.}，日本語の「ら」や連記）を手動で管理すると，
表記揺れや誤記の原因となります．

\textsf{authcite} は，Bib\TeX スタイルファイル（\texttt{authcite.bst}）と連携し，
引用キーから自動的に著者情報を取得して統一された書式で出力します．

\section{インストールと準備}
\subsection{必要ファイル}
\begin{itemize}
  \item \texttt{authcite.sty}: \LaTeX パッケージ本体
  \item \texttt{authcite.bst}: Bib\TeX スタイルファイル（標準 \texttt{unsrt.bst} ベース）
\end{itemize}
これらのファイルを執筆対象の \texttt{.tex} ファイルと同じディレクトリ，
または \TeX システムの検索パス（\texttt{TEXINPUTS} / \texttt{BSTINPUTS}）に配置してください．

\subsection{文書での指定}
\begin{verbatim}
\documentclass{...}
\usepackage{authcite}

\begin{document}
...
\bibliographystyle{authcite}
\bibliography{your_bib_file}
\end{document}
\end{verbatim}

\section{使い方と出力例}

\subsection{基本コマンド}
本文中では \verb|\authcite{key}| を使用します．
\begin{itemize}
  \item \textbf{単著 (1人)}: \authcite{hosseini2020}（コード: \verb|\authcite{hosseini2020}|）
  \item \textbf{2人共著 (既定)}: \authcite{du2021}（コード: \verb|\authcite{du2021}|）
  \item \textbf{3人以上}: \authcite{smith2022}（コード: \verb|\authcite{smith2022}|）
  \item \textbf{others 含む}: \authcite{lee2023}（コード: \verb|\authcite{lee2023}|）
  \item \textbf{複合姓}: \authcite{dimonda2025}（コード: \verb|\authcite{dimonda2025}|）
  \item \textbf{著者なし}: \authcite{manual2024}（コード: \verb|\authcite{manual2024}|）
\end{itemize}

\subsection{2人共著の連記 (\texttt{two=both})}
2人共著の場合，学術機関や学会の執筆要領（情報処理学会など）によっては「〜ら」と省略せず，
両名を連記することが推奨されます．
オプション \texttt{two=both} を指定すると，2名とも表示されます：
\begin{itemize}
  \item 欧文2名連記: \authcite[two=both]{du2021}（コード: \verb|\authcite[two=both]{du2021}|）
  \item 全体の既定値を連記に設定したい場合は，プリアンブルで \verb|\authcitesetup{two=both}| を指定します．
\end{itemize}

\subsection{英語モード (\texttt{lang=en})}
英語論文，または英語での引用箇所では \texttt{lang=en} を指定します：
\begin{itemize}
  \item 単著 (en): \authcite[lang=en]{hosseini2020}（コード: \verb|\authcite[lang=en]{hosseini2020}|）
  \item 2名 et al. (en): \authcite[lang=en]{du2021}（コード: \verb|\authcite[lang=en]{du2021}|）
  \item 2名 連記 (en): \authcite[lang=en,two=both]{du2021}（コード: \verb|\authcite[lang=en,two=both]{du2021}|）
  \item 3名 (en): \authcite[lang=en]{smith2022}（コード: \verb|\authcite[lang=en]{smith2022}|）
\end{itemize}

\section{オプション仕様}
\begin{table}[htbp]
\centering
\caption{設定オプション一覧}
\begin{tabular}{llll}
\hline
キー & 設定値 & 既定値 & 説明 \\
\hline
\texttt{lang} & \texttt{ja} / \texttt{en} & \texttt{ja} & 助詞・表記言語（ja:「ら」「と」／en:「et al.」「and」） \\
\texttt{two}  & \texttt{etal} / \texttt{both} & \texttt{etal} & 2名共著の表記（etal: 1人目＋複数語／both: 2名連記） \\
\hline
\end{tabular}
\end{table}

\section{コンパイルの流れ}
\textsf{authcite} は \textsf{natbib} と同様，Bib\TeX 出力（\texttt{.bbl}）と \texttt{.aux} を介して著者名を受け渡します．
そのため，文書を完全に組版するには以下の順序でコンパイルしてください：
\begin{enumerate}
  \item \texttt{platex} (または \texttt{uplatex})：\texttt{.aux} に引用キーを書き出し
  \item \texttt{pbibtex} (または \texttt{upbibtex})：\texttt{.bbl} に \verb|\authciteentry| を生成
  \item \texttt{platex}：\texttt{.bbl} を読み込み，\texttt{.aux} に著者情報を登録
  \item \texttt{platex}：\texttt{.aux} から著者情報を反映して最終組版
\end{enumerate}
※ \texttt{latexmk} 等のビルドツールを使用している場合は自動的に解決されます．

\bibliographystyle{authcite}
\bibliography{\jobname}

\end{document}

